% META: NONE.1

\section{Характеристика объекта исследования}

\subsection{САПР КОМПАС-3D как целевая платформа}
Объектом исследования является процесс параметрического моделирования в системе автоматизированного проектирования \textbf{КОМПАС-3D} (разработчик — АО «АСКОН»). КОМПАС-3D — одна из ведущих российских САПР, широко применяемая в машиностроении, приборостроении и строительстве. Система поддерживает гибридное параметрическое моделирование, позволяя создавать трёхмерные модели деталей и сборок с полной историей построения. Дерево построения в КОМПАС-3D содержит такие операции, как выдавливание, вырезание, вращение, кинематические операции, скругления, фаски, массивы и др. Каждая операция имеет набор численных параметров и ссылки на геометрические объекты (эскизы, грани, рёбра).

Однако при импорте моделей из нейтральных форматов (STEP, IGES) или при получении геометрии от сторонних разработчиков история построения утрачивается, и пользователь получает «немую» модель — статичное B-Rep тело. Цель настоящей работы — разработать модуль, который по такой «немой» геометрии восстанавливает последовательность операций, приближённую к исходной.

\subsection{Типовые детали для экспериментов}
Для обучения и тестирования использовались реальные машиностроительные детали, предоставленные Заказчиком. Всего в сыром датасете насчитывалось 82~418 моделей, охватывающих широкий спектр сложности — от простых пластин с отверстиями до корпусных деталей с большим числом формообразующих элементов. 

Примеры типовых деталей (см. Приложение Б):
\begin{itemize}
    \item \textbf{Фланец} — плоская деталь с отверстиями по окружности, создаваемыми операцией \texttt{cutExtrusion} (вырезание выдавливанием) и часто содержащая скругления (\texttt{fillet}) на внешних и внутренних рёбрах.
    \item \textbf{Вал ступенчатый} — тело вращения, однако в нашем датасете многие валы моделировались последовательными выдавливаниями (\texttt{bossExtrusion}) с последующими фасками (\texttt{chamfer}) на торцах.
    \item \textbf{Корпусная деталь} — содержит все четыре исследуемые операции: базовое выдавливание для основания, вырезания для полостей и окон, скругления для обеспечения технологичности и фаски для снятия острых кромок.
\end{itemize}

В процессе фильтрации были отобраны модели, содержащие хотя бы одну из четырёх ключевых операций: \texttt{bossExtrusion} (выдавливание), \texttt{cutExtrusion} (вырезание выдавливанием), \texttt{fillet} (скругление), \texttt{chamfer} (фаска). Эти операции составляют основу большинства машиностроительных деталей и обладают хорошо выраженными геометрическими признаками.

\subsection{Типовые операции и статистика датасета}

Выбор именно четырёх операций обусловлен несколькими факторами. Во-первых, согласно анализу частотности в сыром датасете (см. таблицу~\ref{tab:ops-stats}), операции выдавливания и вырезания встречаются более чем в 85\% моделей, скругления — примерно в 40\%, фаски — в 12\%. Во-вторых, эти операции обладают наиболее чёткими геометрическими следами на поверхности, что критически важно для обучения нейросетевой сегментации. В-третьих, они являются базовыми для подавляющего большинства CAD-систем, и их автоматическое распознавание даёт наибольшую практическую пользу для упрощения моделей.

\subsubsection*{Описание операций}

\begin{itemize}
    \item \texttt{bossExtrusion} (выдавливание): добавление объёма путём перемещения плоского эскиза в направлении, перпендикулярном его плоскости. Геометрические следы — плоские грани, параллельные плоскости эскиза, и боковые поверхности, ортогональные к ним. Сложность распознавания — средняя, так как выдавливание может быть как «прилипающим» к существующей геометрии, так и независимым.
    
    \item \texttt{cutExtrusion} (вырезание выдавливанием): удаление объёма по тому же принципу. Геометрический признак — «вырезанная» полость или отверстие, часто сквозное или заданной глубины. Особенность: вырезания могут перекрываться последующими операциями, поэтому их сегментация требует анализа отрицательной геометрии.
    
    \item \texttt{fillet} (скругление): создание плавного перехода между двумя пересекающимися поверхностями путём замены острого ребра цилиндрической или тороидальной поверхностью. Геометрические признаки: плавно сопряжённые поверхности, постоянный или переменный радиус, характерное изменение кривизны. Скругления наиболее легко распознаются нейросетью благодаря гладкости и высокой локальной симметрии.
    
    \item \texttt{chamfer} (фаска): снятие острого ребра под заданным углом, создание плоской наклонной грани вместо ребра. Геометрические признаки: плоская наклонная грань, углы 45° или другие заданные значения. Сложность распознавания выше, чем у скруглений, так как фаска плоская и может быть спутана с другими плоскими гранями, не являющимися результатом отдельной операции.
\end{itemize}

\subsubsection*{Статистика по датасету после фильтрации}

В таблице~\ref{tab:ops-stats} приведены итоговые количества моделей, содержащих каждую операцию, после очистки и приведения к рабочему набору (21~464 модели). Данные получены путём анализа JSON-файлов истории построения (см. раздел «Методика получения статистики» ниже).

\begin{table}[h]
\centering
\caption{Распределение операций в рабочем датасете}
\label{tab:ops-stats}
\begin{tabular}{|l|c|c|c|}
\hline
\textbf{Операция} & \textbf{Кол-во моделей} & \textbf{Доля от общего числа} & \textbf{Среднее число операций на модель} \\
\hline
bossExtrusion & 12\,847 & 59.9\% & 2.4 \\
cutExtrusion & 11\,932 & 55.6\% & 3.1 \\
fillet & 8\,294 & 38.6\% & 1.7 \\
chamfer & 2\,576 & 12.0\% & 1.2 \\
\hline
\end{tabular}
\end{table}

Виден сильный дисбаланс: \texttt{chamfer} представлена примерно в 5 раз реже, чем \texttt{bossExtrusion}. Это потребовало применения методов аугментации данных (небольшие вращения, масштабирование, добавление шума к вершинам) для операции фаски, чтобы избежать переобучения.

Для каждой операции была сформирована своя обучающая выборка, причём размеры выборок различались в соответствии с таблицей. Например, для \texttt{chamfer} было отобрано 2\,576 положительных примеров, и к ним добавлено такое же количество отрицательных (моделей без фаски) для бинарной классификации. Для сегментации использовались все модели, содержащие операцию, с автоматической разметкой рёбер (см. раздел 2.3.1).
